\subsection{EXPONENTIAL}

Let \(\mathrm{E}(x)\) be the expansion of an exponential mapping of G as an integral series about 0. Let \(\mathrm{L}(x)\) be the expansion of the inverse mapping of an injective exponential mapping of G as an integral series about 0. Since the tangent mapping at 0 to any exponential mapping is the identity of \(\mathrm{L}(\mathrm{G})\) , \(\mathrm{E}(x) \equiv x \mod \deg 2\) and \(\mathrm{L}(x) \equiv x \mod \deg 2\) . Since \(\mathrm{E}(\mathrm{L}(x)) = \mathrm{L}(\mathrm{E}(x))\) for \(x\) sufficiently close to 0, the formal power series E and L are such that \(\mathrm{E}(\mathrm{L}(\mathrm{X})) = \mathrm{L}(\mathrm{E}(\mathrm{X})) = \mathrm{X}\) . An analogous argument shows that

\[
\mathrm{E} (t \mathrm{X}) = (\mathrm{E} (\mathrm{X})) ^ {[ t ]}, \quad \mathrm{L} (\mathrm{X} ^ {[ t ]}) = t \mathrm{L} (\mathrm{X})
\]

for \(t \in K\) .

PROPOSITION 3.

\begin{enumerate}
\def\labelenumi{(\arabic{enumi})}
\setcounter{enumi}{26}
\tightlist
\item
\end{enumerate}

\[
\mathrm{L} = \sum_ {p = 1} ^ {\infty} \frac {(- 1) ^ {p - 1}}{p} \psi_ {p}\tag{27}
\]

\[
\mathrm{E} = \sum_ {p = 1} ^ {\infty} \frac {1}{p !} \psi_ {p, p}\tag{28}
\]

(Recall that \(\psi_{p,p}\) is the homogeneous component of \(\psi_{p}\) of degree \(p\) .)

\[
\operatorname{E} (t x) = (\operatorname{E} (x)) ^ {[ t ]}
\]

or, by (24),

\[
\mathrm{E} (t x) = \sum_ {m \geqslant 0} \sum_ {1 \leqslant r \leqslant m} t ^ {r} \phi_ {r, m} (\mathrm{E} (x)).
\]

The two sides are formal power series in \(t\) and \(x\) . Equating the terms of first degree in \(t\) , we obtain

\[
x = \sum_ {m \geqslant 0} \phi_ {1, m} (\mathrm{E} (x)).\tag{29}
\]

Now the inversion of a system of formal power series, when it is possible, is possible in only one way (Algebra, Chapter IV, § 6, Corollary to Proposition 8). Then

\[
\begin{array}{r l} \mathrm{L} (x) & = \sum_ {m \geq 0} \phi_ {1, m} (x) \quad \text { by   (29) } \\ & = \sum_ {p, m} \frac {(- 1) ^ {p}}{p} \psi_ {p, m} (x) \quad \text { by   (25) } \\ & = \sum_ {p} \frac {(- 1) ^ {p}}{p} \psi_ {p} (x), \end{array}
\]

whence (i). Similarly, for \(t \neq 0\) ,

\[
\begin{array}{l} \mathrm{L} (t x) = t \mathrm{L} ((t x) ^ {[ t ^ {- 1} ]}) \\ = t \mathrm{L} \left(\sum_ {m \geqslant 0} \sum_ {1 \leqslant r \leqslant m} t ^ {m - r} \phi_ {r, m} (x)\right) \quad \text { by } (24). \end{array}
\]

Equating the terms of first degree in t, we obtain

\[
x = \mathrm{L} \Bigl (\sum_ {m \geqslant 0} \phi_ {m, m} (x) \Bigr)
\]

whence

\[
\begin{array}{r l} \mathrm{E} (x) & = \sum_ {m \geqslant 0} \phi_ {m, m} (x) \\ & = \sum_ {m \geqslant 0} \frac {1}{m !} \psi_ {m, m} (x) \quad \text { by   (26) }. \end{array}
\]

PROPOSITION 4. For the chart of G used to be canonical, it is necessary and sufficient that \(\psi_j = 0\) for \(j \geqslant 2\) .

This is sufficient by Proposition 3. Suppose that the chart is canonical and that \(\psi_{i} = 0\) for \(2 \leqslant i < n\) . Then \(nx = x^{[n]} = \sum_{i=0}^{n} \binom{n}{i} \psi_{i}(x) = nx + \psi_{n}(x)\) , whence \(\psi_{n} = 0\) . Hence \(\psi_{j} = 0\) for \(j \geqslant 2\) by induction on \(j\) .

